% =====================================================================
%  §5 — Human–AI Annotation Workflow and Evaluation
% =====================================================================

\section{Human--AI Annotation Workflow and Evaluation}
\label{sec:workflow}

We designed and piloted an annotation workflow in which every annotation is
\emph{auditable}. Each assertion is anchored to the specific text that
supports it; each dimension value is defensible under the schema's
activation conditions; and each judgment that a downstream user might
contest carries an explicit reviewer or AI flag.

The same commitment applies to AI assistance. When an AI system drafts an
annotation, its uncertainties and assumptions are recorded for curator
review rather than silently absorbed into the annotation.
\Cref{sec:workflow-flags} reports the practical effect of that review. The
manual and AI-drafted tracks have different epistemic roles, primary and
secondary knowledge work, respectively, and use distinct flag vocabularies
to record provenance.

We treat the pilot as a \emph{feasibility study}, not as a benchmark.
Six papers and a single annotator do not support claims about
inter-annotator agreement or generalizability. The pilot shows that the
workflow's epistemic commitments can be implemented in practice and that
the schema exposes, rather than hides, its points of strain.

\begin{figure}[tbp]
  \centering
  \begin{tikzpicture}[
      box/.style={draw, rounded corners, align=center, inner sep=3pt,
                  text width=25mm, minimum height=9mm, font=\footnotesize},
      arr/.style={-{Latex[length=2mm]}, semithick}]
    \node[box, fill=black!5] (src) {Source publication (PDF)};
    \node[box, fill=blue!6, below left=12mm and 22mm of src] (m1) {Curator highlights \& callouts};
    \node[box, fill=blue!6, below=6mm of m1] (m2) {Structured JSON record};
    \node[box, fill=blue!6, below=6mm of m2] (m3) {YAML annotation\\(curator ground truth)};
    \node[box, fill=orange!8, below right=12mm and 22mm of src] (a1) {AI draft\\(annotation skill)};
    \node[box, fill=orange!8, below=6mm of a1] (a2) {Curator review\\(review skill)};
    \node[box, fill=orange!8, below=6mm of a2] (a3) {Reviewed YAML annotation};
    \node[box, fill=black!5] (val) at ($(m3)!0.5!(a3) - (0,1.3cm)$) {SHACL validation\\(conditional activation)};
    \draw[arr] (src) -- (m1);
    \draw[arr] (src) -- (a1);
    \draw[arr] (m1) -- (m2);
    \draw[arr] (m2) -- (m3);
    \draw[arr] (a1) -- (a2);
    \draw[arr] (a2) -- (a3);
    \draw[arr] (m3) -- (val);
    \draw[arr] (a3) -- (val);
  \end{tikzpicture}
  \caption{The two annotation tracks. Curator-authored annotations
  (left) form the ground truth. AI-drafted annotations (right) are
  produced under an annotation skill from the source PDF, the schema, and
  the manual annotations as exemplars, then reviewed by a curator under a
  review skill. Both tracks are validated against the same SHACL shapes.
  In every case, the curator, not the model, is the authority for the
  recorded annotation.}
  \label{fig:workflow}
\end{figure}

\subsection{Protocol and Coverage}
\label{sec:workflow-protocol}

Annotations are created through two tracks (\cref{fig:workflow}).

\paragraph{Manual track.}
Four of the six publications, Jossin, Davis, Nelson, and Gupta, were
annotated manually by Vladimir Seplyarskiy, a researcher in human and
population genetics. The curator read each paper, marked relevant
passages in the PDF with highlights and callouts, exported those
annotations with a helper script to a structured JSON record, and mapped
the result, together with a per-paper summary row, to a YAML annotation
conforming to the schema.

When this conversion required normalization, the change was recorded
inline as a \texttt{normalization\_note} with
\texttt{type = ai\_normalization}. For example, the free-text
``localisation and protein expression'' was mapped to the enumeration
values \dimval{Localization} and \dimval{Existence}. This record allows
later audits to compare the original and normalized values. The
manual-track protocol is described in more detail in Supplementary
Note~\ref*{supp:sec:supp-protocol-suite}.

\paragraph{AI-drafted track.}
For two publications, Duerr and Inouye, an AI assistant drafted the
annotations directly from the PDFs under an executable annotation
skill~\cite{anthropic_skills}. The assistant received the schema
(\texttt{dimensions.md}), the ground-truth annotations for the four
manual papers as exemplars, and the source PDFs. It produced YAML
annotations in the same format as the manual annotations.

When the assistant judged a mapping to be uncertain, or concluded that a
judgment call required human review, it recorded the issue inline as
either \texttt{ai\_uncertainty} or \texttt{ai\_assumption}. One of the
authors then reviewed the AI-drafted annotations assertion by assertion.
The reviewer has expertise in genomic data curation and the evidence
model, rather than in genetics domain science. The review assessed source
fidelity and schema conformance, not the underlying science.
The four manual annotations are treated as domain-expert ground truth for
their own papers and as exemplars for AI drafting. The executable skills
and review protocol are described in Supplementary
Notes~\ref*{supp:sec:supp-protocol-suite}
and~\ref*{supp:sec:supp-skills}.

\paragraph{Source anchoring.}
In both tracks, every assertion, that is, the \texttt{assertion} property
of a \GeneticEvidence{} item, carries a \texttt{source\_span}. The span
identifies a specific page and includes a short key phrase, typically fewer
than fifteen words, from the source text. The phrase is long enough for a
reader to locate the passage by search while remaining short enough to
respect fair-use norms for copyrighted material.

Source anchoring is mandatory: an annotation with any assertion lacking a
source span is rejected at validation time. Across the six-paper pilot, all
95 assertions are source-anchored. This is a property of the workflow's
validation requirement, not an observed outcome.

\paragraph{Coverage.}
Dimension coverage across the 28 \GeneticEvidence{} items is summarized in
Supplementary Note~\ref*{supp:sec:supp-coverage-note}
(Table~\ref*{supp:tab:supp-coverage}) and in the repository's coverage
report.\footnote{\url{https://github.com/ForomePlatform/genetic-evidence-model/blob/master/annotations/coverage.md}}

\subsection{Flags and Extension Promotion}
\label{sec:workflow-flags}
\label{sec:workflow-promotion}

The two tracks use small flag vocabularies to make their epistemic
provenance explicit.

\paragraph{Manual track.}
The manual track uses three reviewer flags. A
\texttt{reviewer\_disagreement} records a clear factual correction, such
as the organism-set correction in Davis~2011. A
\texttt{reviewer\_suggestion} records an optional addition whose inclusion
depends on curator preference, such as HPO terms for patient phenotypes in
Davis, Nelson, and Gupta. A \texttt{reviewer\_query} records a case in
which the mapping to the schema is unclear and should not be silently
corrected. For example, the Nelson summary assigned
\texttt{specificity\_of\_phenotype = SNV}, which does not fit the
dimension's semantics.

\paragraph{AI-drafted track.}
Because there is no prior ground truth with which to disagree, the
AI-drafted track uses two flags. An \texttt{ai\_uncertainty} records a
mapping that the assistant judged unreliable and referred for a human
decision, such as a borderline assignment of Mode of Inheritance. An
\texttt{ai\_assumption} records a judgment made without strong textual
support, such as whether cited mouse-model evidence should be represented
as a separate \GeneticEvidence{} item or treated as background context.

Across the six-paper pilot, manual-track annotations carried five queries,
six suggestions, and one disagreement, whereas AI-drafted annotations
carried six uncertainties and six assumptions. Of the nine AI-drafted items
from Duerr and Inouye, none was rejected: one was approved as drafted,
three were approved with non-blocking flags, and five were edited. The
later Inouye draft required fewer edits than the earlier Duerr draft, but
the papers differ in genre and protocol version; this comparison is
therefore suggestive rather than controlled. The edits mainly affected
dimension assignment and credibility calibration rather than the factual
content of assertions. The reviewer also added a small number of
higher-level items, such as a paper's overall thesis or a cited
translational study. Per-paper counts and detailed rationales are given in
Supplementary Note~\ref*{supp:sec:supp-ai-review}.

\paragraph{Extension promotion.}
The workflow adopts a conservative rule for evolving the schema. A
representational gap observed in one paper is recorded only in that
paper's \texttt{candidate\_extensions} block. A gap observed
independently in a second paper is promoted to a first-class dimension or
enumeration value.

In the six-paper pilot, this rule promoted two extensions. \emph{Phenotype
Scale} was first proposed in Jossin and then reused in Davis, Nelson, and
later papers. \emph{Variant Ascertainment} was first proposed in Davis and
then reused in Nelson and Duerr. Both now appear in the vocabulary
(\cref{tab:dimensions-core,tab:dimensions-cond}): Phenotype Scale among
the always-required dimensions and Variant Ascertainment among the
conditional dimensions, where it is required when Target Type is
\dimval{Variant}.

Fifteen other candidates remain unpromoted, pending a second independent
use or an explicit curator decision. The rule also prevents premature
additions. For example, a proposed polarity dimension for negative
assertions in Duerr was withdrawn when review confirmed that the existing
assertion predicate already represents both absence and presence.
